\section{The \texttt{MCPL\_input\_once} McStas Component}
Source-like component that reads neutron state parameters from a MCPL-file one time.

\subsection*{Identification}
\begin{itemize}
  \item \textbf{Author:} Gregory S Tucker
  \item \textbf{Origin:} European Spallation Source ERIC
  \item \textbf{Date:} Sep 2024, reduced to an INHERIT of MCPL\_input Oct 2026
\end{itemize}

\subsection*{Description}
Source-like component that reads neutron state parameters from a MCPL-file one time.

Since Oct 2026 MCPL\_input\_once is MCPL\_input under its old name: everything, including the parameters, is inherited from MCPL\_input, which now splits the file between MPI processes the way MCPL\_input\_once did. With the default repeat\_count=1, every event in the file is used exactly once. See MCPL\_input for the full documentation.

MCPL is short for Monte Carlo Particle List, and is a format for sharing events between e.g. MCNP(X), Geant4 and McStas.

When used with MPI, the file contents are shared between workers with each accessing approximately (\#events in the file) / (\#MPI nodes)

Example: MCPL\_input\_once(filename=voutput,v\_smear=0.1,pos\_smear=0.001,dir\_smear=0.01)

\subsection*{Input parameters}
Parameters in \textbf{boldface} are required; the others are optional.

\begin{longtable}{p{0.22\textwidth}p{0.12\textwidth}p{0.46\textwidth}p{0.14\textwidth}}
\toprule
\textbf{Name} & \textbf{Unit} & \textbf{Description} & \textbf{Default} \\
\midrule
\endhead
\bottomrule
\end{longtable}

\subsection*{Links}
\begin{itemize}
  \item Component source code found in file \texttt{MCPL\_input\_once.comp}.
\end{itemize}
\IfFileExists{obsolete/MCPL_input_once_static.tex}{\subsubsection*{Input parameters}
\texttt{MCPL\_input\_once} is defined as
\texttt{DEFINE COMPONENT MCPL\_input\_once INHERIT MCPL\_input} and takes
exactly the input parameters of \texttt{MCPL\_input} (see the Misc chapter),
including \texttt{repeat\_count} (default 1). With the default settings each
event in the MCPL file is used exactly once, also when running with MPI.
Use \texttt{MCPL\_input} in new instruments.
}{}